We will begin by defining some notation. Given a distribution $\p(t)$ and an
integrable function $\phi(t)$, we will use the expectation operator as a formal
shorthand for the integral
%
%\begin{align*}
$\expect{\p(t)}{\phi(t)} = \int \phi(t) \p(dt)$.
%\end{align*}
%
Any quantities not explicitly part of the distribution beneath the expectation
operator should be thought of as fixed. Covariances and variances are treated
similarly, i.e.\,
%
%\begin{align*}
%
$\cov{\p(t)}{\phi(t), \psi(t)} =
    \expect{\p(t)}{\phi(t) \psi(t)^T} -
        \expect{\p(t)}{\phi(t)} \expect{\p(t)}{\psi(t)}^T$.
%
%\end{align*}
%
We will denote sample covariances over an index $n \in \{1,
\ldots, N \} = [N]$ as $\covhat{n \in [N]}{\cdot, \cdot}$.
% %
% \begin{align*}
% %
% \covhat{n \in [N]}{s_n, t_n} :={}& \frac{1}{N}
%     \sumn s_n t_n - \left( \meann s_n \right) \left( \meann t_n \right)^T.
% %
% \end{align*}

We denote partial derivatives with respect to $\theta$ with bracketed
subscripts.  For example, for a function $\psi(\theta)$ taking values in
$\rdom{D_\psi}$, let $\psigrad{k}(\theta)$ denote the array of $k$-th order
partial derivatives with respect to $\theta$, i.e.,
%
%\begin{align*}
%
$
\psigrad{k}(\theta) :=
    \fracat{\partial^k \psi(\theta)}{\partial\theta^k}{\theta}.
$
%
%\end{align*}
%
If $\psi$ depends on any variables other than $\theta$, those variables will be
taken as fixed in the partial differentiation.
%
In general, $\psigrad{k}(\theta)$ is an array with $D_\psi \thetadim^k$ entries, and we
will introduce special notation as needed when summing over the array's entries.
But we will mostly be able to avoid such special notation by interpreting the
following special cases in standard matrix notation:
%
% \begin{itemize}
% %
% \item When $D_\psi = 1$ and $k = 1$, $\psigrad{1}(\theta)$ is taken to be a
% $\thetadim$-column vector.
% %
% \item When $D_\psi > 1$ and $k = 1$, $\psigrad{1}(\theta)$ is taken to be a
% $D_\psi \times \thetadim$ matrix.
% %
% \item When $D_\psi = 1$ and $k=2$, $\psigrad{2}(\theta)$ is taken to be a $\thetadim
% \times \thetadim$ matrix.
% %
% \item When $k = 0$, $\psigrad{0}(\theta)$ is just $\psi(\theta)$.
% %
% \end{itemize}
%
when $D_\psi = 1$ and $k = 1$, $\psigrad{1}(\theta)$ is taken to be a
$\thetadim$-column vector; when $D_\psi > 1$ and $k = 1$, $\psigrad{1}(\theta)$
is taken to be a $D_\psi \times \thetadim$ matrix; when $D_\psi = 1$ and $k=2$,
$\psigrad{2}(\theta)$ is taken to be a $\thetadim \times \thetadim$ matrix; and
when $k = 0$, $\psigrad{0}(\theta)$ is just $\psi(\theta)$.


% Define likelihood and related quantities
Our proposal and existing methods for estimating $\gcovtrue$ will depend on some key
quantities to which we give symbols in \defref{loglik_def}.  In the left hand
column of \defref{loglik_def} we have finite-sample quantities, and on the right
hand side we have the corresponding population quantities. In \secref{bayes_clt}
below we will state precise conditions under which all the quantities of
\defref{loglik_def} are well-defined.  We will use $\xn$ to denote a generic
single datapoint, in contrast to specific observed datapoints, $\x_n$.
%
\begin{defn}\deflabel{loglik_def}
\setlength{\abovedisplayskip}{0pt}
\setlength{\abovedisplayshortskip}{0pt}
\begin{align*}
&\textrm{Finite-sample quantity} && \textrm{Population quantity}\\
%
\likhat(\theta) :={}&
   \frac{1}{N} \sumn \ell(\x_n \vert \theta) +
   \frac{1}{N} \logprior(\theta)
&
\quad \lik(\theta) :={}& \expect{\fdist(\xn)}{\ell(\xn \vert \theta)}
\\
\thetahat :={}& \argmax_{\theta \in \thetadom} \likhat(\theta)
&
\thetatrue :={}& \argmax_{\theta \in \thetadom} \lik(\theta)
\\
\scorecovhat :={}&
   \frac{1}{N} \sumn \ellgrad{1}(\x_n \vert \thetahat)
                     \ellgrad{1}(\x_n \vert \thetahat)^T
&
\scorecov :={}&
    \cov{\fdist(\xn)}{\ellgrad{1}(\xn \vert \thetatrue)}
\\
% \infohat :={}& -\fracat{\partial^2 \likhat(\theta)}
%                      {\partial\theta \partial\theta^T}{\thetahat}
\infohat :={}& -\likhatk{2}(\thetahat)
&
\info :={}& - \expect{\fdist(\xn)}{\ellgrad{2}(\xn \vert \thetatrue)} \qedhere
%
\end{align*}
\end{defn}
%

% We are studying fixed frequentist properties.
We will take the population quantities in the right column of
\defref{loglik_def} to be fixed (though, in general, unknown) and study the
sampling behavior of quantities in the left column. That is, we will study the
asymptotic behavior of formal Bayesian posterior quantities as a function of a
fixed data generating process, following much of the literature on the
frequentist behavior of Bayesian estimators (e.g., \citet{van:2000:asymptotic}).
% Note that the numeric vector $\thetatrue$ is
% conceptually distinct from the random variable $\theta$ of the Bayesian
% posterior, despite also existing in $\thetadom$.
%
Our notation for stochastic convergence reflects these assumptions.  In
particular, order notation $O(\cdot)$ applies as $N \rightarrow \infty$, and
probabilistic order notation $O_p(\cdot)$ denotes order in probability over
$\fdist(\xn)$.  Tilde order notation $\ordlogp{\cdot}$ neglects factors that are
sub-polynomial in $N$.

Let $\norm{\cdot}_1$, $\norm{\cdot}_2$, and $\norm{\cdot}_\infty$ denote the
usual vector norms.  When applied to an array of derivatives, the norms refer to
the vector norm applied to the vectorized version of the array.



% The notation
% $\normdist(\theta | \mu, \Sigma)$ will denote a normal distribution on the
% random variable $\theta$ with mean $\mu$ and covariance $\Sigma$.
